% Lösungen Potenzen, Zehnerpotenzen und Quadratwurzeln (zusammenbau v0.4, Heft)
\einheitenkopf{Potenzen, Zehnerpotenzen und Quadratwurzeln}

\erg{18}{a) $2, 4, 8, 16, 32, 64, 128, 256, 512$: neun Faktoren, also $x = 9$ \quad b) $2, 4, 8, 16, 32, 64, 128$: sieben Faktoren, also $x = 7$ \quad c) $4, 16, 64$: drei Faktoren, also $x = 3$ \quad d) $4, 16, 64, 256, 1\,024$: fünf Faktoren, also $x = 5$ \quad e) $9^3 = 729$, $3^5 = 243$, $2^9 = 512$ – die größte ist $9^3$ \quad f) $6^4 = 1\,296$, $4^5 = 1\,024$, $3^6 = 729$ – die größte ist $6^4$ \quad g) $5^{-3} = \frac{1}{125} = 0{,}008 < 0{,}01$ \quad h) $2^{-4} = \frac{1}{16} = 0{,}0625 > 0{,}05$}
\erg{19}{a) $630\,000$ – das Komma wandert fünf Stellen nach rechts \quad b) $2\,070\,000$ – das Komma wandert sechs Stellen nach rechts \quad c) $9\,100\,000\,000$ – das Komma wandert neun Stellen nach rechts \quad d) $5$ – das Komma wandert fünf Stellen \quad e) $5$ – das Komma wandert fünf Stellen \quad f) $0{,}00034$ – das Komma wandert vier Stellen nach links \quad g) $0{,}000068$ – das Komma wandert fünf Stellen nach links \quad h) $-4 \cdot 10^{3}$, denn $-4 \cdot 10^{3} = -4\,000$ liegt auf der Zahlengerade am weitesten links \quad i) $-7 \cdot 10^{5}$, denn $-7 \cdot 10^{5} = -700\,000$ liegt auf der Zahlengerade am weitesten links}
\erg{20}{a) $\sqrt{6{,}25} = 2{,}5$, die Seite ist $2{,}5$ m lang \quad b) $\sqrt{1\,600} = 40$, die Seite ist $40$ m lang \quad c) $\sqrt{441} = 21$, die Seite ist $21$ cm lang \quad d) $\sqrt{(-7)^2} = 7$, denn erst das Quadrat: $(-7)^2 = 49$, und die Wurzel ist nie negativ \quad e) $\sqrt{(-13)^2} = 13$, denn erst das Quadrat: $(-13)^2 = 169$, und die Wurzel ist nie negativ \quad f) $\frac{7}{3} > \frac{9}{4}$, denn $\frac{7}{3} \approx 2{,}33$ und $\frac{9}{4} = 2{,}25$; $\sqrt{5} \approx 2{,}24$ ist kleiner \quad g) $\sqrt{3} < \frac{7}{4}$, denn $\sqrt{3} \approx 1{,}73$ und $\frac{7}{4} = 1{,}75$ \quad h) $-0{,}26 < -\frac{1}{4} < 1{,}7 < \sqrt{3}$, denn $-\frac{1}{4} = -0{,}25$ und $\sqrt{3} \approx 1{,}73$ \quad i) $-0{,}62 < -\frac{3}{5} < 3{,}3 < \sqrt{11}$, denn $-\frac{3}{5} = -0{,}6$ und $\sqrt{11} \approx 3{,}32$}
